\documentclass[aspectratio=169, 10pt]{beamer}
\usetheme{metropolis}

\usepackage{tikz}
\usetikzlibrary{calc, arrows.meta, shapes.geometric}

% ============================================================
% PRESENTATION META INFORMATION
% ============================================================

\title{Engineering Curves: Construction of Epicycloid}
\subtitle{Rolling Circle Method with Tangent \& Normal}
\author{Gajavada Sanjeevkumar}
\date{}

\setbeamertemplate{frame footer}
{\scriptsize Gajavada Sanjeevkumar \quad|\quad Engineering Curves: Epicycloid}

\setbeamertemplate{page number in head/foot}[totalframenumber]

% ============================================================
% TIKZ STYLES
% ============================================================

\tikzset{
  fixed circle/.style={
    thick,
    blue!80!black
  },

  generating circle/.style={
    thick,
    blue!55!black
  },

  directrix line/.style={
    ultra thick,
    magenta!90!black
  },

  construction line/.style={
    thin,
    blue!60!black
  },

  rolling line/.style={
    thin,
    dashed,
    cyan!70!black
  },

  arc line/.style={
    thin,
    dashed,
    red!85!black
  },

  main curve/.style={
    ultra thick,
    magenta!90!black
  },

  tangent line/.style={
    thick,
    green!50!black
  },

  normal line/.style={
    thick,
    orange!80!black
  },

  axis line/.style={
    thin,
    dash pattern=on 6pt off 2pt on 1.5pt off 2pt,
    black!75
  },

  label line/.style={
    stealth-,
    thin,
    black!70
  }
}

% ============================================================
% LEGEND CARD
% ============================================================

\newcommand{\drawLegendCard}[1]{
  \node[
    draw=gray!50,
    fill=white,
    rounded corners=4pt,
    inner sep=4pt,
    anchor=north east
  ] at (8.0,5.2) {
    \scriptsize
    \begin{tabular}{l@{\,$\rightarrow$\,}l}
      #1
    \end{tabular}
  };
}

% ============================================================
% DOCUMENT
% ============================================================

\begin{document}

% ============================================================
% TITLE SLIDE
% ============================================================

\begin{frame}[plain]
  \titlepage
\end{frame}


% ============================================================
% PROBLEM STATEMENT
% ============================================================

\begin{frame}{Problem Statement}

\vfill

\begin{block}{Question}

Draw an \textbf{epicycloid} generated by a circle of radius

\[
r=25\text{ mm}
\]

rolling externally without slipping around a fixed circle of radius

\[
R=50\text{ mm}.
\]

Construct the curve for \textbf{one complete revolution of the generating-circle centre}.

Also construct the \textbf{tangent} and \textbf{normal} at a selected point $P$ on the epicycloid.

\end{block}

\vfill

\end{frame}


% ============================================================
% STEP 1
% FIXED CIRCLE + GENERATING CIRCLE
% ============================================================

\begin{frame}{Step 1: Draw the Fixed Circle and Generating Circle}

\centering

\begin{tikzpicture}[scale=0.75, >=stealth]

  \useasboundingbox (-5,-5) rectangle (8,5.5);

  \drawLegendCard{
    \textbf{$O$} & Centre of Fixed Circle\\
    \textbf{$R$} & Fixed Circle Radius\\
    \textbf{$r$} & Generating Circle Radius
  }

  % Fixed circle
  \coordinate (O) at (0,0);

  \fill[red!80!black] (O) circle (3pt)
    node[below left,font=\small] {$O$};

  \draw[fixed circle] (O) circle (3);

  % Initial generating circle
  \coordinate (C0) at (4.5,0);

  \fill[red!80!black] (C0) circle (3pt)
    node[below right,font=\small] {$C_0$};

  \draw[generating circle] (C0) circle (1.5);

  % Common contact point
  \coordinate (P0) at (6,0);

  \fill[blue!80!black] (P0) circle (3pt)
    node[above right,font=\small] {$P_0$};

  % Radius R
  \draw[<->,thin,red!80!black]
    (0,0) -- (3,0)
    node[midway,below=3pt,font=\tiny] {$R=50$ mm};

  % Radius r
  \draw[<->,thin,cyan!80!black]
    (4.5,0) -- (6,0)
    node[midway,above=4pt,font=\tiny] {$r=25$ mm};

  % Centre line
  \draw[axis line]
    (-3.5,0) -- (6.5,0);

  % External tangency indication
  \draw[construction line,dashed]
    (O) -- (P0);

  \node[below=18pt,font=\small]
    at (3,0)
    {Initial external contact};

\end{tikzpicture}

\end{frame}


% ============================================================
% STEP 2
% CENTRE LOCUS
% ============================================================

\begin{frame}{Step 2: Draw the Locus of the Generating-Circle Centre}

\centering

\begin{tikzpicture}[scale=0.75, >=stealth]

  \useasboundingbox (-5,-5) rectangle (8,5.5);

  \drawLegendCard{
    \textbf{$O$} & Fixed Centre\\
    \textbf{$R+r$} & Centre-Locus Radius\\
    \textbf{$C$} & Generating-Circle Centre
  }

  \coordinate (O) at (0,0);

  % Fixed circle
  \draw[fixed circle] (O) circle (3);

  % Centre locus
  \draw[directrix line,dashed] (O) circle (4.5);

  % Initial centre
  \coordinate (C0) at (4.5,0);

  \fill[red!80!black] (O) circle (3pt)
    node[below left,font=\small] {$O$};

  \fill[red!80!black] (C0) circle (3pt)
    node[right,font=\small] {$C_0$};

  % Radius R+r
  \draw[<->,thin,orange!80!black]
    (0,0) -- (4.5,0)
    node[midway,below=4pt,font=\tiny] {$R+r=75$ mm};

  % Generating circle
  \draw[generating circle] (C0) circle (1.5);

  % Axis
  \draw[axis line]
    (-5,0) -- (7,0);

  \node[below=20pt,font=\small]
    at (2.25,0)
    {Locus of centre $C$: radius $R+r$};

\end{tikzpicture}

\end{frame}


% ============================================================
% STEP 3
% DIVIDE FIXED CIRCLE
% ============================================================

\begin{frame}{Step 3: Divide the Fixed Circle into 12 Equal Parts}

\centering

\begin{tikzpicture}[scale=0.75, >=stealth]

  \useasboundingbox (-5,-5) rectangle (8,5.5);

  \drawLegendCard{
    \textbf{$O$} & Fixed Centre\\
    \textbf{$1\dots12$} & Twelve Equal Divisions\\
    \textbf{$C$-locus} & Centre Locus
  }

  \coordinate (O) at (0,0);

  % Fixed circle
  \draw[fixed circle] (O) circle (3);

  % Centre locus
  \draw[directrix line,dashed] (O) circle (4.5);

  % Radial divisions
  \foreach \a [count=\i] in
    {0,30,60,90,120,150,180,210,240,270,300,330} {

      \draw[construction line,dashed]
        (O) -- ++(\a:3);

      \fill[blue!80!black]
        ({3*cos(\a)},{3*sin(\a)}) circle (1.5pt);

  }

  % Division labels
  \node[right,font=\tiny] at (3,0) {1};
  \node[above right,font=\tiny] at (2.598,1.5) {2};
  \node[above right,font=\tiny] at (1.5,2.598) {3};
  \node[above,font=\tiny] at (0,3) {4};
  \node[above left,font=\tiny] at (-1.5,2.598) {5};
  \node[above left,font=\tiny] at (-2.598,1.5) {6};
  \node[left,font=\tiny] at (-3,0) {7};
  \node[below left,font=\tiny] at (-2.598,-1.5) {8};
  \node[below left,font=\tiny] at (-1.5,-2.598) {9};
  \node[below,font=\tiny] at (0,-3) {10};
  \node[below right,font=\tiny] at (1.5,-2.598) {11};
  \node[below right,font=\tiny] at (2.598,-1.5) {12};

  \fill[red!80!black] (O) circle (3pt)
    node[below left,font=\small] {$O$};

\end{tikzpicture}

\end{frame}


% ============================================================
% STEP 4
% LOCATE CENTRES
% ============================================================

\begin{frame}{Step 4: Locate Centre Positions $C_1\dots C_{12}$}

\centering

\begin{tikzpicture}[scale=0.75, >=stealth]

  \useasboundingbox (-5,-5) rectangle (8,5.5);

  \drawLegendCard{
    \textbf{$C_0$} & Initial Centre\\
    \textbf{$C_1\dots C_{12}$} & Successive Centres\\
    \textbf{$R+r$} & Centre-Locus Radius
  }

  \coordinate (O) at (0,0);

  \draw[fixed circle] (O) circle (3);

  \draw[directrix line,dashed]
    (O) circle (4.5);

  % Centres
  \foreach \a [count=\i] in
    {0,30,60,90,120,150,180,210,240,270,300,330} {

      \coordinate (C\i) at
        ({4.5*cos(\a)},{4.5*sin(\a)});

      \fill[red!80!black]
        (C\i) circle (2pt);

      \node[
        font=\tiny,
        red!80!black
      ] at
        ({4.5*cos(\a)},{4.5*sin(\a)+0.22})
        {$C_{\i}$};

      \draw[rolling line]
        (O) -- (C\i);
  }

  \fill[red!80!black] (O) circle (3pt)
    node[below left,font=\small] {$O$};

\end{tikzpicture}

\end{frame}


% ============================================================
% STEP 5
% ROLLING CIRCLE POSITIONS
% ============================================================

\begin{frame}{Step 5: Draw Successive Positions of the Generating Circle}

\centering

\begin{tikzpicture}[scale=0.75, >=stealth]

  \useasboundingbox (-6,-6) rectangle (7,6);

  \drawLegendCard{
    \textbf{$C_1\dots C_{12}$} & Successive Centres\\
    \textbf{Circle} & Generating Circle\\
    \textbf{Rotation} & Rolling Without Slip
  }

  \coordinate (O) at (0,0);

  % Fixed circle
  \draw[fixed circle] (O) circle (3);

  % Centre locus
  \draw[directrix line,dashed] (O) circle (4.5);

  % Successive centre positions
  \foreach \a [count=\i] in
    {0,30,60,90,120,150,180,210,240,270,300,330} {

      \coordinate (C\i) at
        ({4.5*cos(\a)},{4.5*sin(\a)});

      \fill[red!80!black]
        (C\i) circle (1.5pt);

      \draw[generating circle,opacity=0.35]
        (C\i) circle (1.5);

  }

  % Highlight three important positions
  \coordinate (C1) at (4.5,0);
  \coordinate (C4) at (0,4.5);
  \coordinate (C7) at (-4.5,0);

  \draw[generating circle]
    (C1) circle (1.5);

  \draw[generating circle]
    (C4) circle (1.5);

  \draw[generating circle]
    (C7) circle (1.5);

  \fill[red!80!black] (O) circle (3pt)
    node[below left,font=\small] {$O$};

  \node[below,font=\small]
    at (0,-5.2)
    {The generating circle rolls externally around the fixed circle};

\end{tikzpicture}

\end{frame}


% ============================================================
% STEP 6
% ROLLING ANGLE RELATION
% ============================================================

\begin{frame}{Step 6: Determine the Rolling Rotation}

\begin{block}{No-Slip Condition}

When the generating circle rolls externally around the fixed circle,

\[
\text{distance travelled by centre}
=
\text{arc rolled on generating circle}.
\]

For centre rotation $\theta$,

\[
(R+r)\theta=r\phi
\]

where $\phi$ is the rotation of the generating circle.

Hence,

\[
\boxed{
\phi=\frac{R+r}{r}\theta
}
\]

For

\[
R=50\text{ mm},\qquad r=25\text{ mm},
\]

we obtain

\[
\boxed{
\phi=3\theta
}
\]

Therefore, during one complete revolution of the centre,

\[
\theta=360^\circ
\]

and the generating circle rotates through

\[
\boxed{
\phi=1080^\circ
}
\]

or three complete rotations.

\end{block}

\end{frame}


% ============================================================
% STEP 7
% GENERATE POINTS
% ============================================================

\begin{frame}{Step 7: Locate Points $P_1\dots P_{12}$}

\centering

\begin{tikzpicture}[scale=0.75, >=stealth]

  \useasboundingbox (-6,-6) rectangle (7,6);

  \drawLegendCard{
    \textbf{$C_1\dots C_{12}$} & Centre Positions\\
    \textbf{$P_1\dots P_{12}$} & Epicycloid Points\\
    \textbf{$r$} & Generating Radius
  }

  \coordinate (O) at (0,0);

  % Fixed circle
  \draw[fixed circle] (O) circle (3);

  % Centre locus
  \draw[directrix line,dashed] (O) circle (4.5);

  % Generate points
  \foreach \a [count=\i] in
    {0,30,60,90,120,150,180,210,240,270,300,330} {

      \pgfmathsetmacro{\px}
        {4.5*cos(\a) + 1.5*cos(-3*\a)}

      \pgfmathsetmacro{\py}
        {4.5*sin(\a) + 1.5*sin(-3*\a)}

      \coordinate (P\i) at (\px,\py);

      \fill[blue!80!black]
        (P\i) circle (2pt);

      \node[
        font=\tiny,
        blue!80!black
      ] at (\px,\py+0.25)
        {$P_{\i}$};

      % centre
      \coordinate (C\i) at
        ({4.5*cos(\a)},{4.5*sin(\a)});

      \fill[red!80!black]
        (C\i) circle (1.5pt);

      % construction radius
      \draw[arc line]
        (C\i) -- (P\i);
  }

  \fill[red!80!black]
    (O) circle (3pt)
    node[below left,font=\small] {$O$};

\end{tikzpicture}

\end{frame}


% ============================================================
% STEP 8
% COMPLETE EPICYCLOID
% ============================================================

\begin{frame}{Step 8: Smoothly Join $P_0\dots P_{12}$ to Complete the Epicycloid}

\centering

\begin{tikzpicture}[scale=0.75, >=stealth]

  \useasboundingbox (-6,-6) rectangle (7,6);

  \drawLegendCard{
    \textbf{Fixed Circle} & Directing Circle\\
    \textbf{Centre Locus} & $R+r$\\
    \textbf{Curve} & Epicycloid Profile
  }

  \coordinate (O) at (0,0);

  % Fixed circle
  \draw[fixed circle] (O) circle (3);

  % Centre locus
  \draw[directrix line,dashed]
    (O) circle (4.5);

  % Epicycloid
  %
  % x = (R+r) cos(theta) - r cos(3theta)
  % y = (R+r) sin(theta) - r sin(3theta)
  %
  \draw[
    main curve,
    samples=250,
    domain=0:360,
    variable=\t
  ]
    plot (
      {4.5*cos(\t) - 1.5*cos(3*\t)},
      {4.5*sin(\t) - 1.5*sin(3*\t)}
    );

  % Selected construction points
  \foreach \a [count=\i] in
    {0,30,60,90,120,150,180,210,240,270,300,330} {

      \pgfmathsetmacro{\px}
        {4.5*cos(\a) + 1.5*cos(-3*\a)}

      \pgfmathsetmacro{\py}
        {4.5*sin(\a) + 1.5*sin(-3*\a)}

      \fill[blue!80!black]
        (\px,\py) circle (1.5pt);
  }

  \fill[red!80!black]
    (O) circle (3pt)
    node[below left,font=\small] {$O$};

  \node[
    below,
    font=\small
  ] at (0,-5.5)
  {Three-cusped epicycloid};

\end{tikzpicture}

\end{frame}


% ============================================================
% STEP 9
% MARK POINT P
% ============================================================

\begin{frame}{Step 9: Mark Point $P$ on the Epicycloid}

\centering

\begin{tikzpicture}[scale=0.75, >=stealth]

  \useasboundingbox (-6,-6) rectangle (7,6);

  \drawLegendCard{
    \textbf{$P$} & Selected Point\\
    \textbf{$C_p$} & Centre Corresponding to $P$\\
    \textbf{$M$} & Instantaneous Contact Point
  }

  \coordinate (O) at (0,0);

  % Fixed circle
  \draw[fixed circle] (O) circle (3);

  % Centre locus
  \draw[directrix line,dashed]
    (O) circle (4.5);

  % Epicycloid
  \draw[
    main curve,
    samples=250,
    domain=0:360,
    variable=\t
  ]
    plot (
      {4.5*cos(\t) - 1.5*cos(3*\t)},
      {4.5*sin(\t) - 1.5*sin(3*\t)}
    );

  % Choose theta = 30 degrees
  \pgfmathsetmacro{\thetaP}{30}

  \coordinate (Cp) at
    ({4.5*cos(\thetaP)},{4.5*sin(\thetaP)});

  \coordinate (M) at
    ({3*cos(\thetaP)},{3*sin(\thetaP)});

  \pgfmathsetmacro{\px}
    {4.5*cos(\thetaP)+1.5*cos(-3*\thetaP)}

  \pgfmathsetmacro{\py}
    {4.5*sin(\thetaP)+1.5*sin(-3*\thetaP)}

  \coordinate (P) at (\px,\py);

  % Centre
  \fill[red!80!black]
    (Cp) circle (3pt)
    node[above right,font=\small] {$C_p$};

  % Point P
  \fill[black]
    (P) circle (3.5pt)
    node[above right,font=\small] {$P$};

  % Contact point M
  \fill[orange!90!black]
    (M) circle (3pt)
    node[below left,font=\small] {$M$};

  % CP
  \draw[construction line,dashed]
    (Cp) -- (P)
    node[midway,above,font=\tiny] {$r$};

  % OM
  \draw[construction line,dashed]
    (O) -- (M);

  % CM
  \draw[construction line,dashed]
    (Cp) -- (M);

  \fill[red!80!black]
    (O) circle (3pt)
    node[below left,font=\small] {$O$};

\end{tikzpicture}

\end{frame}


% ============================================================
% STEP 10
% NORMAL
% ============================================================

\begin{frame}{Step 10: Draw Normal $N$-$N'$ Through $P$ and $M$}

\centering

\begin{tikzpicture}[scale=0.75, >=stealth]

  \useasboundingbox (-6,-6) rectangle (7,6);

  \drawLegendCard{
    \textbf{$M$} & Instantaneous Contact Point\\
    \textbf{$P$} & Point on Epicycloid\\
    \textbf{$N$-$N'$} & Normal
  }

  \coordinate (O) at (0,0);

  % Fixed circle
  \draw[fixed circle] (O) circle (3);

  % Centre locus
  \draw[directrix line,dashed]
    (O) circle (4.5);

  % Epicycloid
  \draw[
    main curve,
    samples=250,
    domain=0:360,
    variable=\t
  ]
    plot (
      {4.5*cos(\t)-1.5*cos(3*\t)},
      {4.5*sin(\t)-1.5*sin(3*\t)}
    );

  % Selected point theta = 30
  \pgfmathsetmacro{\thetaP}{30}

  \coordinate (Cp) at
    ({4.5*cos(\thetaP)},{4.5*sin(\thetaP)});

  \coordinate (M) at
    ({3*cos(\thetaP)},{3*sin(\thetaP)});

  \pgfmathsetmacro{\px}
    {4.5*cos(\thetaP)+1.5*cos(-3*\thetaP)}

  \pgfmathsetmacro{\py}
    {4.5*sin(\thetaP)+1.5*sin(-3*\thetaP)}

  \coordinate (P) at (\px,\py);

  % Normal direction vector P-M
  \pgfmathsetmacro{\dx}{\px-3*cos(\thetaP)}
  \pgfmathsetmacro{\dy}{\py-3*sin(\thetaP)}

  % Normal extended line
  \draw[
    normal line
  ]
    ($(P)!-1.8!(M)$)
      node[below left,font=\small] {$N$}
    --
    ($(M)!2.0!(P)$)
      node[above right,font=\small] {$N'$};

  % M-P construction
  \draw[construction line,dashed]
    (M) -- (P);

  % Points
  \fill[orange!90!black]
    (M) circle (3pt)
    node[below left,font=\small] {$M$};

  \fill[black]
    (P) circle (3.5pt)
    node[above right,font=\small] {$P$};

  \fill[red!80!black]
    (O) circle (3pt)
    node[below left,font=\small] {$O$};

  \fill[red!80!black]
    (Cp) circle (2pt);

\end{tikzpicture}

\end{frame}


% ============================================================
% STEP 11
% TANGENT
% ============================================================

\begin{frame}{Step 11: Draw Tangent $T$-$T'$ Perpendicular to the Normal}

\centering

\begin{tikzpicture}[scale=0.75, >=stealth]

  \useasboundingbox (-6,-6) rectangle (7,6);

  \drawLegendCard{
    \textbf{$N$-$N'$} & Normal\\
    \textbf{$T$-$T'$} & Tangent\\
    \textbf{$P$} & Point of Contact on Curve
  }

  \coordinate (O) at (0,0);

  % Fixed circle
  \draw[fixed circle] (O) circle (3);

  % Centre locus
  \draw[directrix line,dashed]
    (O) circle (4.5);

  % Epicycloid
  \draw[
    main curve,
    samples=250,
    domain=0:360,
    variable=\t
  ]
    plot (
      {4.5*cos(\t)-1.5*cos(3*\t)},
      {4.5*sin(\t)-1.5*sin(3*\t)}
    );

  % Selected point
  \pgfmathsetmacro{\thetaP}{30}

  \coordinate (M) at
    ({3*cos(\thetaP)},{3*sin(\thetaP)});

  \pgfmathsetmacro{\px}
    {4.5*cos(\thetaP)+1.5*cos(-3*\thetaP)}

  \pgfmathsetmacro{\py}
    {4.5*sin(\thetaP)+1.5*sin(-3*\thetaP)}

  \coordinate (P) at (\px,\py);

  % Normal vector
  \pgfmathsetmacro{\nx}{\px-3*cos(\thetaP)}
  \pgfmathsetmacro{\ny}{\py-3*sin(\thetaP)}

  % Normal length
  \pgfmathsetmacro{\nlen}{sqrt(\nx*\nx+\ny*\ny)}

  % Unit normal
  \pgfmathsetmacro{\ux}{\nx/\nlen}
  \pgfmathsetmacro{\uy}{\ny/\nlen}

  % Unit tangent = (-uy, ux)
  \pgfmathsetmacro{\tx}{-\uy}
  \pgfmathsetmacro{\ty}{\ux}

  % Normal
  \draw[
    normal line
  ]
    ({\px-3*\ux},{\py-3*\uy})
      node[below left,font=\small] {$N$}
    --
    ({\px+3*\ux},{\py+3*\uy})
      node[above right,font=\small] {$N'$};

  % Tangent
  \draw[
    tangent line
  ]
    ({\px-3*\tx},{\py-3*\ty})
      node[below left,font=\small] {$T$}
    --
    ({\px+3*\tx},{\py+3*\ty})
      node[above right,font=\small] {$T'$};

  % Right angle marker
  \coordinate (Q1) at
    ({\px-0.35*\tx},{\py-0.35*\ty});

  \coordinate (Q2) at
    ({\px-0.35*\tx+0.35*\ux},
     {\py-0.35*\ty+0.35*\uy});

  \coordinate (Q3) at
    ({\px+0.35*\ux},
     {\py+0.35*\uy});

  \draw[thin,black]
    (Q1) -- (Q2) -- (Q3);

  % Points
  \fill[black]
    (P) circle (3.5pt)
    node[above right,font=\small] {$P$};

  \fill[orange!90!black]
    (M) circle (3pt)
    node[below left,font=\small] {$M$};

  \fill[red!80!black]
    (O) circle (3pt)
    node[below left,font=\small] {$O$};

\end{tikzpicture}

\end{frame}


% ============================================================
% STEP 12
% FINAL DRAWING
% ============================================================

\begin{frame}{Step 12: Final Completed Epicycloid}

\centering

\begin{tikzpicture}[scale=0.75, >=stealth]

  \useasboundingbox (-6,-6) rectangle (7,6);

  \drawLegendCard{
    \textbf{Curve} & Completed Epicycloid\\
    \textbf{$N$-$N'$} & Normal\\
    \textbf{$T$-$T'$} & Tangent
  }

  \coordinate (O) at (0,0);

  % ==========================================================
  % FIXED CIRCLE
  % ==========================================================

  \draw[fixed circle]
    (O) circle (3);

  % ==========================================================
  % CENTRE LOCUS
  % ==========================================================

  \draw[directrix line,dashed]
    (O) circle (4.5);

  % ==========================================================
  % EPICYCLOID
  % ==========================================================

  \draw[
    main curve,
    samples=300,
    domain=0:360,
    variable=\t
  ]
    plot (
      {4.5*cos(\t)-1.5*cos(3*\t)},
      {4.5*sin(\t)-1.5*sin(3*\t)}
    );

  % ==========================================================
  % CONSTRUCTION POINTS
  % ==========================================================

  \foreach \a in
    {0,30,60,90,120,150,180,210,240,270,300,330} {

      \pgfmathsetmacro{\px}
        {4.5*cos(\a)+1.5*cos(-3*\a)}

      \pgfmathsetmacro{\py}
        {4.5*sin(\a)+1.5*sin(-3*\a)}

      \fill[blue!80!black]
        (\px,\py) circle (1.5pt);
  }

  % ==========================================================
  % SELECTED POINT P
  % ==========================================================

  \pgfmathsetmacro{\thetaP}{30}

  \coordinate (M) at
    ({3*cos(\thetaP)},{3*sin(\thetaP)});

  \pgfmathsetmacro{\px}
    {4.5*cos(\thetaP)+1.5*cos(-3*\thetaP)}

  \pgfmathsetmacro{\py}
    {4.5*sin(\thetaP)+1.5*sin(-3*\thetaP)}

  \coordinate (P) at (\px,\py);

  % Normal vector
  \pgfmathsetmacro{\nx}{\px-3*cos(\thetaP)}
  \pgfmathsetmacro{\ny}{\py-3*sin(\thetaP)}

  \pgfmathsetmacro{\nlen}
    {sqrt(\nx*\nx+\ny*\ny)}

  \pgfmathsetmacro{\ux}{\nx/\nlen}
  \pgfmathsetmacro{\uy}{\ny/\nlen}

  % Tangent vector
  \pgfmathsetmacro{\tx}{-\uy}
  \pgfmathsetmacro{\ty}{\ux}

  % ==========================================================
  % NORMAL
  % ==========================================================

  \draw[
    normal line
  ]
    ({\px-3*\ux},{\py-3*\uy})
      node[below left,font=\small] {$N$}
    --
    ({\px+3*\ux},{\py+3*\uy})
      node[above right,font=\small] {$N'$};

  % ==========================================================
  % TANGENT
  % ==========================================================

  \draw[
    tangent line
  ]
    ({\px-3*\tx},{\py-3*\ty})
      node[below left,font=\small] {$T$}
    --
    ({\px+3*\tx},{\py+3*\ty})
      node[above right,font=\small] {$T'$};

  % ==========================================================
  % P AND M
  % ==========================================================

  \fill[black]
    (P) circle (3.5pt)
    node[above right,font=\small] {$P$};

  \fill[orange!90!black]
    (M) circle (3pt)
    node[below left,font=\small] {$M$};

  % Normal construction segment
  \draw[construction line,dashed]
    (M) -- (P);

  % Centre
  \fill[red!80!black]
    (O) circle (3pt)
    node[below left,font=\small] {$O$};

  % ==========================================================
  % RIGHT ANGLE
  % ==========================================================

  \coordinate (Q1) at
    ({\px-0.35*\tx},{\py-0.35*\ty});

  \coordinate (Q2) at
    ({\px-0.35*\tx+0.35*\ux},
     {\py-0.35*\ty+0.35*\uy});

  \coordinate (Q3) at
    ({\px+0.35*\ux},
     {\py+0.35*\uy});

  \draw[thin,black]
    (Q1) -- (Q2) -- (Q3);

\end{tikzpicture}

\end{frame}


% ============================================================
% END
% ============================================================

\end{document}
